\documentclass[a4paper,12pt,oneside]{book}

% pakiety
\usepackage{polski}
%\usepackage[utf8]{inputenc} % nie można używać as per 
			     % https://tex.stackexchange.com/a/480069/177956
\usepackage{fancyhdr} % nagłówki i stopki
\usepackage{indentfirst} % WAŻNE, MA BYĆ!
\usepackage{graphicx} % to do wstawiania rysunków
\usepackage{amsmath} % to do dodatkowych symboli, przydatne
\usepackage[pdftex,
            left=1in,right=1in,
            top=1in,bottom=1in]{geometry} % marginesy
\usepackage{amssymb} % to też do dodatkowych symboli, też przydatne
\usepackage{pdfpages}
\usepackage{lipsum}
\usepackage{multirow}
\usepackage{listings}
\usepackage{caption}
\usepackage{booktabs}
\usepackage{subcaption}
\usepackage{xcolor}
\graphicspath{ {./img/} }
\DeclareCaptionType{code}[Listing][Spis listingów] 

\definecolor{codegreen}{rgb}{0,0.6,0}
\definecolor{codegray}{rgb}{0.5,0.5,0.5}
\definecolor{codepurple}{rgb}{0.58,0,0.82}
\definecolor{backcolour}{rgb}{0.95,0.95,0.92}

\lstset{
	backgroundcolor=\color{backcolour},   
	commentstyle=\color{codegreen},
	keywordstyle=\color{magenta},
	numberstyle=\tiny\color{codegray},
	stringstyle=\color{codepurple},
	basicstyle=\footnotesize,
	breakatwhitespace=false,         
	breaklines=true,                 
	captionpos=b,                    
	keepspaces=true,                 
	numbers=left,                    
	numbersep=5pt,                  
	showspaces=false,                
	showstringspaces=false,
	showtabs=false,                  
	tabsize=2,
	float=h
}

% definicje nagłówków i stopek
\pagestyle{fancy}
\renewcommand{\chaptermark}[1]{\markboth{#1}{}}
\renewcommand{\sectionmark}[1]{\markright{\thesection\ #1}}
% komenda która sprawi że nie numerowany chapter (chapter*) jest dodany do TOC
\newcommand\chap[1]{%
  \chapter*{#1}%
  \addcontentsline{toc}{chapter}{#1}}
\fancyhf{}
\fancyhead[LE,RO]{\footnotesize\bfseries\thepage}
\fancyhead[LO]{\footnotesize\rightmark}
\fancyhead[RE]{\footnotesize\leftmark}
\renewcommand{\headrulewidth}{0.5pt}
\renewcommand{\footrulewidth}{0pt}
\addtolength{\headheight}{1.5pt}
\fancypagestyle{plain}{\fancyhead{}\cfoot{\footnotesize\thepage}\renewcommand{\headrulewidth}{0pt}}


% interlinia
\linespread{1.25}


% treść
\begin{document}
% strona tytułowa
\sloppy
\thispagestyle{empty}
\includepdf{strona_tytulowa}
\newpage{}

\thispagestyle{empty}
\newpage{}

% spis treści
\tableofcontents{}

\chap{Wstęp}
Jedną z większych dziedzin uczenia maszynowego jest uczenie przez wzmacnianie (ang.
\textit{reinforcement learning}) W odróżnieniu od zarówno uczenia nadzorowanego i 
nienadzorowanego nie potrzebujemy w tym przypadku żadnych gotowych danych wejściowych i
wyjściowych. Dzięki temu, że algorytmy uczenia przez wzmacnianie nie mają tego
ograniczenia możemy zastosować je do problemów takich jak gra na giełdzie\cite{
tradingreinforcement}
Jakie rodzaje problemów mogą być rozwiązane przez reinforcement learning?\\
Celem pracy jest zaimplementowanie algorytmu uczenia przez wzmacnianie -- Q-Learning oraz zoptymalizowaniu go tak by po zastosowaniu go do własnoręcznie zaimplementowanej gry był w stanie osiągnąć w niej możliwie najwyższy wynik.



\addcontentsline{toc}{chapter}{Bibliografia}
\bibliographystyle{ieeetr}
\bibliography{praca_licencjacka}

\end{document}
